%%%%%%%%%%%%%%%%%%%%%%%%%%%%%%%%%%%%%%%%%%%%%%%%%%%%%%%%%%%%%%%%%%%%%%%%%%%
\chapter*{Résumé}
\phantomsection
\addcontentsline{toc}{chapter}{Abstract et Résumé}
\thispagestyle{empty}
%%%%%%%%%%%%%%%%%%%%%%%%%%%%%%%%%%%%%%%%%%%%%%%%%%%%%%%%%%%%%%%%%%%%%%%%%%%
Ce projet de fin d’études a été réalisé dans le cadre d’un stage au sein de la société \textbf{Yes Internet}. Il consiste à concevoir et développer une application web permettant la gestion centralisée de plusieurs agences de location de véhicules.

L’objectif principal de cette plateforme est de digitaliser et de faciliter la gestion des activités liées à la location de voitures, notamment la gestion des véhicules, des clients et des contrats de location. Le système intègre également un mécanisme de scoring client permettant d’évaluer la fiabilité des utilisateurs à partir de leur historique de location.

La solution propose aussi un module de gestion des paiements permettant d’enregistrer et de suivre les transactions effectuées par les clients. Des tableaux de bord interactifs sont mis à la disposition des administrateurs et des agences afin de visualiser les informations importantes et d’améliorer le suivi de l’activité.

Sur le plan technique, l’application repose sur une architecture web moderne basée sur le framework Laravel pour le développement du backend, React pour l’interface utilisateur et MySQL pour la gestion de la base de données. Cette architecture garantit une solution performante, évolutive et adaptée aux besoins des agences de location de véhicules.



